% !TeX root = ../main.tex
\chapter{Các phương pháp cải thiện}
\label{chap:improvements}

\section{Thiết kế tuning chung}

Mỗi phương pháp thay đổi một cơ chế regularization so với baseline và được
tuning riêng bằng \texttt{GridSearchCV}, 5 folds, scoring là accuracy. CV chỉ
nhìn training data. Hình~\ref{fig:tuning} cho thấy mean CV accuracy và độ lệch
chuẩn tại từng ứng viên.

\begin{figure}[H]
\centering
\includegraphics[width=\textwidth]{hyperparameter_tuning.png}
\caption{Kết quả 5-fold cross-validation cho ba phương pháp}
\label{fig:tuning}
\end{figure}

\begin{table}[H]
\centering
\caption{Hyperparameter được chọn trên tập training}
\label{tab:selected-parameters}
\begin{tabular}{lll}
\toprule
Phương pháp & Giá trị được chọn & Best CV accuracy \\
\midrule
Giới hạn độ sâu & \texttt{max\_depth=4} & 0,9385 \\
Mẫu tối thiểu tại leaf & \texttt{min\_samples\_leaf=3} & 0,9363 \\
Cost-complexity pruning & \texttt{ccp\_alpha=0.005934} & 0,9341 \\
\bottomrule
\end{tabular}
\end{table}

\section{Improvement 1: giới hạn \texttt{max\_depth}}

\textbf{What.} Notebook thử các độ sâu 2, 3, 4, 5, 6, 7, 8 và 10. Giới hạn
độ sâu ngăn cây tiếp tục chia những nhánh quá đặc thù; nếu giá trị quá nhỏ thì
cây có thể underfit.

\textbf{Result.} CV chọn \texttt{max\_depth=4}. Trên test set, accuracy tăng
từ 91,23\% lên 93,86\% và error rate giảm từ 8,77\% xuống 6,14\%. Cây còn 11
leaf, training accuracy 98,68\%, weighted F1 93,87\% và malignant recall
92,86\%.

\textbf{Why.} Việc loại bỏ ba tầng sâu làm giảm các split phụ thuộc training
nhưng vẫn giữ các quyết định quan trọng ở gần root. Khoảng cách train--test giảm
từ 8,77 xuống 4,82 điểm phần trăm, cho thấy generalization tốt hơn.

\begin{figure}[H]
\centering
\includegraphics[width=\textwidth]{tree_max_depth.png}
\caption{Decision Tree sau khi chọn \texttt{max\_depth=4}}
\label{fig:tree-max-depth}
\end{figure}

\section{Improvement 2: điều chỉnh \texttt{min\_samples\_leaf}}

\textbf{What.} Notebook thử các giá trị 1, 2, 3, 4, 5, 6, 8, 10 và 15. Tham
số buộc mỗi leaf phải có đủ mẫu, hạn chế leaf chỉ đại diện một vài quan sát.

\textbf{Result.} CV chọn \texttt{min\_samples\_leaf=3}. Cây giảm còn depth 6
và 13 leaf; tuy nhiên test accuracy vẫn là 91,23\%, error rate 8,77\%, trong
khi malignant recall giảm từ 92,86\% xuống 90,48\%.

\textbf{Why.} Regularization đã giảm độ phức tạp và training accuracy còn
98,02\%, nhưng ràng buộc này không cải thiện ranh giới quyết định trên test
split. Một vài leaf nhỏ của baseline có thể vẫn chứa thông tin hữu ích; gộp
chúng làm thay đổi hai phía của confusion matrix mà không giảm tổng số lỗi.

\begin{figure}[H]
\centering
\includegraphics[width=\textwidth]{tree_min_samples_leaf.png}
\caption{Decision Tree với \texttt{min\_samples\_leaf=3}}
\label{fig:tree-min-leaf}
\end{figure}

\section{Improvement 3: cost-complexity pruning}

\textbf{What.} Các ứng viên \texttt{ccp\_alpha} được lấy từ pruning path của
training set. Tham số này cân bằng độ phù hợp với dữ liệu và chi phí do thêm
nhánh; CV chọn mức pruning thay vì nhìn test accuracy.

\textbf{Result.} Giá trị \texttt{ccp\_alpha=0.005934} tạo cây depth 3, 6 leaf.
Test accuracy đạt 93,86\%, error rate 6,14\%, training accuracy 97,36\% và
weighted F1 93,84\%. Malignant recall là 90,48\%.

\textbf{Why.} Pruning loại bỏ 13 trong số 19 leaf của baseline và giảm khoảng
cách train--test còn 3,50 điểm phần trăm. Mô hình tổng quát tốt hơn và dễ đọc
hơn rõ rệt. Tuy nhiên, mức cắt này bỏ một nhánh giúp phát hiện thêm một ca
malignant, nên recall malignant thấp hơn mô hình Max Depth.

\begin{figure}[H]
\centering
\includegraphics[width=0.92\textwidth]{tree_pruned.png}
\caption{Decision Tree sau cost-complexity pruning}
\label{fig:tree-pruned}
\end{figure}
